\section{Conclusion}
\label{sec:conclusion}
We presented a parametric, fully labeled staircase point-cloud benchmark together with an apples-to-apples comparison of four classical geometric methods: RANSAC plane fitting, PCA normal analysis, DBSCAN with slope confirmation, and height-histogram banding. All four operate under a unified three-class formulation of treads, risers, and other points. On the steppable tread class, orientation-aware DBSCAN with per-cluster slope confirmation was the strongest method, reaching an F1 of $0.943$ and an IoU of $0.892$, with PCA-normals a close second. These findings rest on a single synthetic staircase under a controlled noise model, so they should be read as relative method rankings pending real-sensor validation. RANSAC was the weakest, primarily because horizontal planes fitted at tread heights absorb riser points and depress precision to $0.714$; the height-histogram, by contrast, attained the highest recall but conflated riser bands with treads. A noise-sensitivity sweep on the PCA-Normals method showed that its performance holds up at realistic jitter but degrades sharply once the height noise approaches half the step height, with tread F1 falling from $0.956$ at $\sigma_z=0.01$\,m to $0.543$ at $\sigma_z=0.07$\,m.

From a safety standpoint, an IoU near $0.89$ is adequate for coarse foothold planning. Yet no method reached unit precision, so residual mislabeling of risers as steppable remains a hazard for stair-following robots and assistive devices, which argues for conservative tread-edge margins and downstream verification. Future work will pursue ensemble fusion of the complementary geometric cues to recover precision, validation on real LiDAR staircases beyond the synthetic generator, and a controlled comparison against learned point-cloud segmenters \cite{qi2016pointnet,sarker2024comprehensive} under low-label regimes, where the reproducibility and interpretability of classical methods matter most.
